\chapter[32]{Nonsingular Solutions on Closed 3-Manifolds}
\label{ch:chow-iv-nonsingular-solutions-closed-three-manifolds}
\dcref{ch32}
\source{chow-et-al-ricci-flow-part-iv:page-0218}

Throughout this chapter all $3$-manifolds are orientable.  The normalized
Ricci flow on a closed $n$-manifold is
\begin{equation*}
 \frac{\partial}{\partial t}g_{ij}
 =-2R_{ij}+\frac{2}{n}r g_{ij},
 \qquad
 r(t)=\frac{\int_{\mathcal M}R_{g(t)}\,d\mu_{g(t)}}
 {\int_{\mathcal M}d\mu_{g(t)}}.
 \tag{32.1}
\end{equation*}
It preserves volume.

\section{Nonsingular solutions and the classification theorem}

\begin{definition}[Nonsingular solution (Definition 32.1)]
\label{def:chow-iv-nonsingular-solution}
\source{chow-et-al-ricci-flow-part-iv:page-0221}
\dcref{ch32:32.1}
\group{Nonsingular Solutions}

A solution $(\mathcal M^n,g(t))$ of the normalized Ricci flow on a closed
manifold is \emph{nonsingular} if it is defined for every
$t\in[0,\infty)$ and there is a constant $C<\infty$ such that
\[
 |\Rm_{g(t)}|\le C
 \qquad\text{on }\mathcal M\times[0,\infty).
\]
\end{definition}

\begin{theorem}[Nonsingular solutions on $3$-manifolds admit geometric decompositions (Theorem 32.2)]
\label{thm:chow-iv-nonsingular-geometric-decomposition}
\source{chow-et-al-ricci-flow-part-iv:page-0222,chow-et-al-ricci-flow-part-iv:page-0223,chow-et-al-ricci-flow-part-iv:page-0224}
\uses{def:chow-iv-nonsingular-solution}
\dcref{ch32:32.2}
\group{Nonsingular Solutions}


If a closed $3$-manifold $\mathcal M$ admits a nonsingular solution of the
normalized Ricci flow, then $\mathcal M$ is diffeomorphic to one of the
following:
\begin{enumerate}
\item a graph manifold;
\item a spherical space form;
\item a flat manifold;
\item a hyperbolic manifold;
\item a union, along incompressible $2$-tori, of finite-volume hyperbolic
manifolds and graph manifolds.
\end{enumerate}
In every case, $\mathcal M$ admits a geometric decomposition in the sense of
Thurston.
\end{theorem}

\begin{definition}[Sequential collapse]
\label{def:chow-iv-sequential-collapse}
\source{chow-et-al-ricci-flow-part-iv:page-0222}
\uses{def:chow-iv-nonsingular-solution}
\dcref{ch32:sequential-collapse}
\group{Nonsingular Solutions}


For a nonsingular solution, put
\[
 \rho(t)=\max_{x\in\mathcal M}\inj_{g(t)}(x).
\]
The solution \emph{sequentially collapses} if there are times
$t_i\to\infty$ such that
\[
 \rho(t_i)^2\max_{\mathcal M\times\{t_i\}}|\Rm|\longrightarrow0.
\]
Otherwise it is \emph{noncollapsed}.  Equivalently, in the noncollapsed
case there is $\delta>0$ such that for every $t\ge0$ some $x_t\in\mathcal M$
satisfies $\inj_{g(t)}(x_t)\ge\delta$.
\end{definition}

\section{The three scalar-curvature cases}

For a normalized Ricci flow let
$R_{\min}(t)=\min_{x\in\mathcal M}R(x,t)$.  Its scalar curvature satisfies
\begin{equation*}
 \partial_tR=\Delta R+2|\Rc|^2-\frac2n rR
 \ge \Delta R+\frac2nR(R-r).
 \tag{32.3}
\end{equation*}
Consequently its lower Dini derivative satisfies
\begin{equation*}
 \frac{d^-}{dt}R_{\min}(t)
 \ge\frac2nR_{\min}(t)(R_{\min}(t)-r(t)).
 \tag{32.4}
\end{equation*}

\begin{lemma}[Monotonicity of $R_{\min}(t)$ (Lemma 32.3)]
\label{lem:chow-iv-minimum-scalar-curvature-monotonicity}
\source{chow-et-al-ricci-flow-part-iv:page-0224}
\dcref{ch32:32.3}
\group{Scalar Curvature Cases}

Let $(\mathcal M^n,g(t))$ be a normalized Ricci flow on a closed manifold.
\begin{enumerate}
\item If $R_{\min}(t_0)>0$, then $R_{\min}(t)>0$ for every $t\ge t_0$.
\item If $R_{\min}(t)\le0$ for all $t$, then $R_{\min}$ is nondecreasing.
If in addition the flow is not Einstein, then it is strictly increasing.
\end{enumerate}
\end{lemma}
\begin{proof}
\sketch
While $R_{\min}>0$, (32.4) gives
\[
 \frac d{dt}\log R_{\min}
 \ge \frac2n(R_{\min}-r)>-\frac2n r.
\]
Since $r$ is bounded on every compact time interval, a positive minimum
cannot reach zero in finite time.  If $R_{\min}\le0$, then
$R_{\min}-r\le0$, so the right-hand side of (32.4) is nonnegative.
If equality occurs at two distinct times, monotonicity makes
$R_{\min}$ constant between them.  The strong maximum principle applied to
(32.3) then makes $R$ spatially constant and forces equality in
$|\Rc|^2\ge R^2/n$; hence $g(t)$ is Einstein.  Thus a non-Einstein flow has
strictly increasing $R_{\min}$.
\end{proof}

\section{Compactness for the normalized Ricci flow}

\begin{definition}[Pointed $C^\infty$ Cheeger--Gromov convergence (Definition 32.4)]
\label{def:chow-iv-pointed-cheeger-gromov-convergence-nrf}
\source{chow-et-al-ricci-flow-part-iv:page-0225}
\dcref{ch32:32.4}
\group{Normalized Ricci Flow Compactness}

A sequence of pointed normalized Ricci flows
$(\mathcal M_i^n,g_i(t),x_i)$ on a common interval $(\alpha,\omega)$
converges to a complete pointed family
$(\mathcal M_\infty^n,g_\infty(t),x_\infty)$ if there are open sets
$U_i\subset\mathcal M_\infty$ exhausting $\mathcal M_\infty$, with
$x_\infty\in U_i$, and diffeomorphisms
\[
 \Phi_i:U_i\longrightarrow V_i=\Phi_i(U_i)\subset\mathcal M_i,
 \qquad \Phi_i(x_\infty)=x_i,
\]
such that $\Phi_i^*(g_i(t)|_{V_i})$ converges to $g_\infty(t)$ in
$C^\infty$, uniformly on compact subsets of
$\mathcal M_\infty\times(\alpha,\omega)$.
\end{definition}

\begin{theorem}[Hamilton's Cheeger--Gromov-type compactness (Theorem 32.5)]
\label{thm:chow-iv-hamilton-compactness-nrf}
\source{chow-et-al-ricci-flow-part-iv:page-0225,chow-et-al-ricci-flow-part-iv:page-0226}
\uses{def:chow-iv-pointed-cheeger-gromov-convergence-nrf}
\dcref{ch32:32.5}
\group{Normalized Ricci Flow Compactness}


Let $(\mathcal M_i^n,g_i(t),x_i)$ be complete pointed normalized Ricci
flows on $(\alpha,\omega)$, where $0\in(\alpha,\omega)$, and suppose that
\[
 |\Rm_{g_i}|\le C_0
 \quad\text{on }\mathcal M_i\times(\alpha,\omega),
 \qquad
 \inj_{g_i(0)}(x_i)\ge\iota_0>0,
\]
with constants independent of $i$.  A subsequence converges in pointed
$C^\infty$ Cheeger--Gromov sense to a complete pointed family
$(\mathcal M_\infty^n,g_\infty(t),x_\infty)$ with bounded curvature that
solves
\begin{equation*}
 \partial_t(g_\infty)_{ij}
 =-2(\Rc_{g_\infty})_{ij}+\frac2n r_\infty(t)(g_\infty)_{ij},
 \qquad
 r_\infty(t)=\lim_{i\to\infty}r_{g_i(t)}.
 \tag{32.5}
\end{equation*}
\end{theorem}
\begin{proof}
\sketch
The curvature bound also bounds every $r_{g_i}$.  On each smaller time
interval, conversion to ordinary Ricci flow followed by Shi's estimates
gives uniform bounds for all spatial curvature derivatives.  Differentiating
the formula for the average scalar curvature shows in turn that all time
derivatives of $r_{g_i}$ are uniformly bounded; each is an average of a
polynomial in $\Rm$ and its spatial derivatives.  Arzela--Ascoli therefore
gives a smooth subsequential limit $r_\infty(t)$.  Hamilton compactness,
using the injectivity-radius hypothesis at time zero, then gives a complete
pointed smooth limit.  Passing to the limit in the normalized flow equation
proves (32.5), and the original curvature bound passes to the limit.
\end{proof}

\begin{example}[Pathologies for averages under pointed convergence]
\label{ex:chow-iv-average-scalar-curvature-pathologies}
\source{chow-et-al-ricci-flow-part-iv:page-0226}
\uses{def:chow-iv-pointed-cheeger-gromov-convergence-nrf}
\dcref{ch32:32.6}
\group{Normalized Ricci Flow Compactness}


For every $n\ge3$, there are closed pointed manifolds whose average scalar
curvatures tend to $+\infty$ while their pointed limit is complete,
noncompact, and has bounded curvature.  For every $n>3$, there is a
bounded-curvature sequence with a complete noncompact pointed limit whose
average scalar curvatures have two distinct subsequential limits.
\end{example}

\begin{remark}[Compactness does not control global averages (Remark 32.7)]
\label{rem:chow-iv-compactness-global-averages}
\source{chow-et-al-ricci-flow-part-iv:page-0226}
\uses{ex:chow-iv-average-scalar-curvature-pathologies}
\dcref{ch32:32.7}
\group{Normalized Ricci Flow Compactness}


Pointed Cheeger--Gromov convergence alone does not force convergence of
average scalar curvatures; global volume can escape through regions that
leave every compact neighborhood of the base point.  Thus a compactness
statement for normalized Ricci flow must retain convergence of the averages
as a separate conclusion obtained from the evolution and derivative bounds.
\end{remark}

\begin{definition}[Limit solution at infinite time]
\label{def:chow-iv-nrf-limit-solution}
\source{chow-et-al-ricci-flow-part-iv:page-0226}
\uses{def:chow-iv-sequential-collapse, thm:chow-iv-hamilton-compactness-nrf}
\dcref{ch32:eq-32.6--32.8}
\group{Normalized Ricci Flow Compactness}


Let $(\mathcal M^3,g(t))$ be a noncollapsed nonsingular normalized Ricci
flow.  Given $t_i\to\infty$, choose $x_i$ with
$\inj_{g(t_i)}(x_i)\ge\delta$ and set
\[
 g_i(t)=g(t_i+t).
\]
Any pointed subsequential limit
$(\mathcal M_\infty^3,g_\infty(t),x_\infty)$, defined for all
$t\in\mathbb R$, is a \emph{limit solution} corresponding to
$\{(x_i,t_i)\}$.  It has bounded curvature and
\[
 \Vol(g_\infty(t))\le\Vol(g(0)).
\]
\end{definition}

\section{The positive and zero cases}

Let $\widetilde g(\widetilde t)$ be the ordinary Ricci flow with
$\widetilde g(0)=g(0)$.  The two flows are related by
\begin{align*}
 g(t)&=\psi(\widetilde t)\widetilde g(\widetilde t),
 &t&=\int_0^{\widetilde t}\psi(\tau)\,d\tau, \tag{32.9}\\
 \psi(\widetilde t)
 &=\exp\left(\frac23\int_0^{\widetilde t}\widetilde r(\tau)\,d\tau\right)
 =\exp\left(\frac23\int_0^t r(\tau)\,d\tau\right). \tag{32.10}
\end{align*}

\begin{lemma}[Criterion for $g(t)$ to have nonnegative curvature (Lemma 32.8)]
\label{lem:chow-iv-limit-nonnegative-curvature-criterion}
\source{chow-et-al-ricci-flow-part-iv:page-0227,chow-et-al-ricci-flow-part-iv:page-0228}
\uses{def:chow-iv-nrf-limit-solution}
\dcref{ch32:32.8}
\group{Positive and Zero Cases}


Let $(\mathcal M^3,g(t))$ be a noncollapsed nonsingular normalized Ricci
flow and let $[0,\widetilde T)$ be the maximal interval for its corresponding
ordinary Ricci flow.  Suppose
\[
 \widetilde t_i\to\widetilde T,
 \qquad
 \widetilde t_i\psi(\widetilde t_i)\to\infty,
 \qquad
 t_i=\int_0^{\widetilde t_i}\psi(\tau)\,d\tau.
 \tag{32.11}
\]
For any base points satisfying the noncollapsing injectivity estimate, every
corresponding limit solution has underlying manifold diffeomorphic to
$\mathcal M$ and is either flat or has positive sectional curvature.
Consequently, $\mathcal M$ is diffeomorphic to a space form of nonnegative
curvature.
\end{lemma}
\begin{proof}
\sketch
Let $\widetilde\nu$ be the smallest curvature-operator eigenvalue of
$\widetilde g$ and put
$C=-\inf_{\mathcal M}\widetilde\nu(\cdot,0)>0$ when this infimum is
negative.  The Hamilton--Ivey estimate gives, wherever
$\widetilde\nu<0$,
\[
 |\widetilde\nu|
 \bigl(\log|\widetilde\nu|+\log(C^{-1}+\widetilde t)-3\bigr)
 \le\widetilde R.
\]
After rescaling by $\psi(\widetilde t_i)$, hypothesis (32.11) makes the
logarithmic factor tend to infinity.  The ratio of scale factors at bounded
translated times stays between fixed positive constants because
\[
 \left|\log\frac{\psi(\widetilde t_i+\widetilde t)}
 {\psi(\widetilde t_i)}\right|
 \le C|\widetilde t|.
\]
Passing to the pointed limit therefore eliminates every negative
curvature-operator eigenvalue, so $g_\infty(t)$ has nonnegative sectional
curvature.

The limit has finite volume.  A complete noncompact manifold with
nonnegative Ricci curvature has infinite volume, so $\mathcal M_\infty$ is
compact; smooth convergence then makes it diffeomorphic to $\mathcal M$.
The classification of closed $3$-dimensional Ricci flows with nonnegative
sectional curvature leaves a flat manifold, a positively curved manifold,
or a finite quotient of a product $S^2\times\mathbb R$.  The product case
cannot be eternal: on a finite cover $S^2\times S^1$, the sphere factor
satisfies
\[
 \partial_th_\infty
 =\left(-R(h_\infty)+\frac23r_\infty\right)h_\infty,
 \qquad
 \frac d{dt}\Area(h_\infty)=-\frac83\pi,
 \tag{32.12}
\]
and hence becomes singular in finite time.  Thus only the flat and
positively curved alternatives remain; in the latter, Hamilton's theorem
identifies the manifold as a spherical space form.
\end{proof}

\begin{remark}[Alternative exclusion of the product case (Remark 32.9)]
\label{rem:chow-iv-exclude-sphere-line-product}
\source{chow-et-al-ricci-flow-part-iv:page-0228}
\uses{lem:chow-iv-limit-nonnegative-curvature-criterion}
\dcref{ch32:32.9}
\group{Positive and Zero Cases}


The $S^2\times\mathbb R$ quotient can alternatively be excluded by applying
Perelman's no-local-collapsing theorem to the eternal limit as
$t\to-\infty$.
\end{remark}

\begin{proposition}[The positive case (Proposition 32.10)]
\label{prop:chow-iv-positive-case-spherical-space-form}
\source{chow-et-al-ricci-flow-part-iv:page-0228,chow-et-al-ricci-flow-part-iv:page-0229}
\uses{lem:chow-iv-minimum-scalar-curvature-monotonicity, lem:chow-iv-limit-nonnegative-curvature-criterion}
\dcref{ch32:32.10}
\group{Positive and Zero Cases}


If a noncollapsed nonsingular normalized Ricci flow
$(\mathcal M^3,g(t))$ satisfies $R_{\min}(t_0)>0$ for some $t_0$, then
$\mathcal M$ is diffeomorphic to a spherical space form.
\end{proposition}
\begin{proof}
\sketch
The corresponding unnormalized Ricci flow has finite maximal time
$\widetilde T$.  Hamilton--Ivey pinching supplies times
$\widetilde t_i\nearrow\widetilde T$ with
$\widetilde R_{\max}(\widetilde t_i)\to\infty$.  Since normalized curvature
is uniformly bounded,
\[
 \widetilde R_{\max}(\widetilde t_i)
 =\psi(\widetilde t_i)R_{\max}(t_i)
 \le C\psi(\widetilde t_i),
\]
so the sequence meets the scaling hypothesis of
Lemma~\ref{lem:chow-iv-limit-nonnegative-curvature-criterion}.
Thus $\mathcal M$ admits a space-form metric of nonnegative curvature.
It also admits the positive-scalar-curvature metric $g(t_0)$, and the
Gromov--Lawson and Schoen--Yau obstruction rules out a flat topology.
Therefore it is a spherical space form.
\end{proof}

\begin{proposition}[The zero case (Proposition 32.11)]
\label{prop:chow-iv-zero-case-flat-metric}
\source{chow-et-al-ricci-flow-part-iv:page-0229,chow-et-al-ricci-flow-part-iv:page-0230}
\uses{lem:chow-iv-limit-nonnegative-curvature-criterion, def:chow-iv-nrf-limit-solution}
\dcref{ch32:32.11}
\group{Positive and Zero Cases}


If a noncollapsed nonsingular normalized Ricci flow
$(\mathcal M^3,g(t))$ on a closed manifold satisfies
\[
 \lim_{t\to\infty}R_{\min}(t)=0,
\]
then $\mathcal M$ admits a flat metric.
\end{proposition}
\begin{proof}
\sketch
Let $\widetilde g$ be the corresponding ordinary Ricci flow.  If its
maximal time $\widetilde T$ is finite, the argument of
Proposition~\ref{prop:chow-iv-positive-case-spherical-space-form} produces a
compact nonnegatively curved limit.  It cannot have positive sectional
curvature, since smooth convergence would make $R_{g(t)}>0$ at some finite
time.  Hence it is flat.

Assume now $\widetilde T=\infty$.  Its volume obeys
\begin{equation*}
 \frac d{d\widetilde t}\Vol(\widetilde g(\widetilde t))
 =-\widetilde r(\widetilde t)\Vol(\widetilde g(\widetilde t)).
 \tag{32.13}
\end{equation*}
There are three possibilities along sequences tending to infinity.

If $\Vol(\widetilde g(\widetilde t_i))\to0$, then
\[
 \Vol(\widetilde g(\widetilde t_i))
 =\psi(\widetilde t_i)^{-3/2}\Vol(g(t_i))
\]
forces $\psi(\widetilde t_i)\to\infty$.  Lemma~\ref{lem:chow-iv-limit-nonnegative-curvature-criterion}
gives a flat or positively curved limit, and the latter is again excluded by
$R_{\min}(t)\to0$.

If $\Vol(\widetilde g(\widetilde t_i))\to\infty$, integration of (32.13)
gives normalized times $t_i\to\infty$ with
$\int_{\mathcal M}R_{g(t_i)}\,d\mu_{g(t_i)}<0$.  A corresponding limit has
nonnegative scalar curvature because $R_{\min}(t_i)\to0$, while
\[
 \int_{\mathcal M}
 \bigl(R_{g(t_i)}-R_{\min}(t_i)\bigr)\,d\mu_{g(t_i)}\longrightarrow0.
\]
Thus the limit has $R\equiv0$.  The normalized scalar-curvature evolution
\[
 \partial_tR=\Delta R+2|\Rc|^2-\frac23r_\infty R
 \tag{32.18}
\]
and the strong maximum principle imply $\Rc\equiv0$ for all earlier times.
In dimension three the limit is flat; its finite volume makes it compact,
so it is diffeomorphic to $\mathcal M$.

In the remaining case the unnormalized volumes stay between two positive
constants.  Equivalently, $\psi$ stays between two positive constants.
Lemma~\ref{lem:chow-iv-limit-nonnegative-curvature-criterion} again gives a
closed flat limit, since positive scalar curvature is impossible.  Each
case therefore equips $\mathcal M$ with a flat metric.
\end{proof}

\section{The negative case}

After rescaling, the negative case is normalized by
\begin{equation*}
 \lim_{t\to\infty}R_{\min}(t)=-6.
 \tag{32.19}
\end{equation*}

\begin{proposition}[Case III: Any limit is hyperbolic (Proposition 32.12)]
\label{prop:chow-iv-negative-case-hyperbolic-limit}
\source{chow-et-al-ricci-flow-part-iv:page-0231,chow-et-al-ricci-flow-part-iv:page-0232}
\uses{lem:chow-iv-minimum-scalar-curvature-monotonicity, def:chow-iv-nrf-limit-solution}
\dcref{ch32:32.12}
\group{Negative Case}


Let $(\mathcal M^3,g(t))$ be a noncollapsed nonsingular normalized Ricci
flow satisfying (32.19).  For any $t_i\to\infty$ and base points satisfying
the noncollapsing injectivity estimate, every corresponding limit solution
is a complete, compact or noncompact, hyperbolic $3$-manifold.  Moreover,
for every $t\in\mathbb R$,
\begin{equation*}
 r_\infty(t)=\lim_{i\to\infty}r_{g_i(t)}=-6,
 \tag{32.20}
\end{equation*}
and $\Vol(g_\infty(t))\le\Vol(g(t))$.  Such a limit is called a
\emph{hyperbolic limit} of $(\mathcal M^3,g(t))$.
\end{proposition}
\begin{proof}
\sketch
Since $R_{\min}(t)\le-6$, (32.4) in dimension three yields
\[
 \frac d{dt}R_{\min}(t)
 \ge4\bigl(r(t)-R_{\min}(t)\bigr)\ge0.
\]
Integration gives
\begin{equation*}
 \int_0^\infty\bigl(r(t)-R_{\min}(t)\bigr)\,dt
 \le -6-R_{\min}(0)<\infty.
 \tag{32.21}
\end{equation*}
The nonnegative integrand has arbitrarily small integral on every translated
unit interval sufficiently far forward.  Passing to a limit solution and
using $R_{\min}(t_i+t)\to-6$ shows that
\[
 \int_t^{t+1}(r_\infty(\tau)+6)\,d\tau=0
\]
for every $t\in\mathbb R$, hence (32.20).

Writing $A=\Vol(g(t))$, which is constant along the normalized flow,
\[
 \int_{\mathcal M}|R-r|\,d\mu
 \le2A(r-R_{\min}).
\]
Together with (32.21), this implies that the integral of $|R-r|$ on every
translated unit space-time cylinder tends to zero.  Smooth pointed
convergence then gives $R_{g_\infty}\equiv-6$.  The scalar-curvature
evolution becomes
\[
 0=2|\Rc_{g_\infty}+2g_\infty|^2,
\]
so
\begin{equation*}
 \Rc_{g_\infty}=-2g_\infty.
 \tag{32.23}
\end{equation*}
In dimension three this is equivalent to sectional curvature $-1$.
Completeness follows from compactness, and lower semicontinuity of volume
under pointed smooth convergence gives
$\Vol(g_\infty(t))\le A$.
\end{proof}
